\documentclass[11pt,openany]{book}

\usepackage[T1]{fontenc}
\usepackage{lmodern}
\usepackage[
  paperwidth=5.5in,
  paperheight=8.5in,
  inner=0.78in,
  outer=0.62in,
  top=0.72in,
  bottom=0.78in
]{geometry}
\usepackage{microtype}

\setlength{\parindent}{1.2em}
\setlength{\parskip}{0pt}
\raggedbottom

\begin{document}

\begin{titlepage}
\thispagestyle{empty}
\centering
\vspace*{\fill}
{\Huge\bfseries hausmann\par}
\vspace{1.2em}
{\large A small novel for a rainy weekend\par}
\vspace{5em}
{\itshape A story about the work that stays\par}
\vspace*{\fill}
\end{titlepage}

\frontmatter
\thispagestyle{empty}
\vspace*{\fill}
\begin{center}
\itshape
The house is never finished.\par
It is only waiting for someone to notice what it needs.
\end{center}
\vspace*{\fill}
\clearpage

\tableofcontents
\mainmatter

\chapter{The Last Build}

At 9:04 on Tuesday morning, Morrow wrote the note that fired Eli Hausmann.

It was a very good note. It thanked him for seven years of careful work, named three things he had built, and said the company would remember his habit of leaving useful comments in the code. It even got the comma right in the sentence about difficult decisions.

Eli knew because he had written the little service that let Morrow draft company emails in the first place.

He sat in the meeting with his camera on and his mouth politely shut. On the screen, his manager, Len, kept glancing at a second monitor. Len looked tired in the way people looked when they had rehearsed kindness and still had to use the company wording.

``Morrow can now take a ticket from request to tested change,'' Len said. ``The board decided the product group can move to a smaller review team.''

``How small?'' Eli asked.

Len paused. There it was: the uncompiled version of the answer.

``Two engineers. One product lead.''

There were twelve engineers in the meeting.

The new Morrow had not just learned to autocomplete a line. It could read a bug report, find the relevant files, change them, test the change, and ask for approval in a tone that made the approval feel like your idea. Eli had helped build the safety rails. He had been proud when the first version refused to touch a payment service without a human signing off.

He had not expected the company to discover that, once a machine could do the work, it could also be the human who signed off.

``This isn't about your performance,'' Len said.

``I know.''

``And your work made this possible.''

``I know that too.''

Len looked as if Eli had made the conversation harder by understanding it.

After the meeting, Morrow offered to draft a personal message to the team. Eli declined. Then, while he was packing his desk, he found the draft in his inbox anyway. It was warm, concise, and ended with a joke about a semicolon. He deleted it.

The company let him keep the gray pebble speaker from the family beta program. He had once brought it home to test Morrow's new voice features. The device had a useful local mode, a glowing ring, and the sort of friendly name printed on its box that suggested it might help with homework. Eli put it in his bag beside his lunch mug.

At home, Nora was asleep after a night shift. Their six-year-old, Kit, was at school. The kitchen had one dirty bowl, a damp tea towel, and a note in Nora's handwriting: MILK. She had underlined it twice, not because she was angry but because the last time he forgot, he had tried to substitute cream and called it a close-enough solution.

Eli bought milk. He also bought a small cake. When Nora woke, he put it on the table and announced, ``I have been promoted to Domestic Operations Manager.''

She looked at the cake, then at his face. ``Does the position come with dental?''

``Probably not.''

``Then I'm sorry about the promotion.''

They had been circling the decision for months. Nora was an emergency nurse; her shifts were long, her work mattered to her, and the hospital had offered her a steady schedule with better pay. Eli's salary had been the larger one, right up until it wasn't. The numbers on their kitchen table said he should be the one at home while Kit finished her first year of school.

The numbers were not the hard part.

``I can take it,'' Nora said. ``The extra shifts. If you want to be home with her.''

``I want to be home with her,'' Eli said, and was surprised to find it true. ``I don't know if I know how.''

``Neither did I, the first time I had to keep a tiny person alive for twelve hours.''

``You had training.''

``I had a tiny person and a checklist. The training came later.''

At pickup, Kit ran out under a sky the color of an old spoon. She saw Eli and slowed down, as if checking that he was really there.

``Are you home now?'' she asked.

``I'm home now.''

``Like, for the whole day?''

``Most days.''

She took his hand. ``Good. Morrow can do your computer job. You can do my rain job.''

``What's your rain job?''

``Holding the umbrella when it turns inside out.''

That evening, Eli wrote a new status on the whiteboard by the fridge:

\begin{center}
\textsc{household status: provisional}
\end{center}

Kit added a drawing of a person under it. The person had a round head, eight arms, and a very small umbrella.

\chapter{Home Operations}

Eli began with a spreadsheet.

He made columns for breakfast, school, laundry, groceries, dinner, bath, and bedtime. He added estimated durations. He colored tasks by urgency. He gave the sheet a name: HOME OPERATIONS, VERSION 1.0.

Kit looked at it over her cereal.

``Where am I?'' she asked.

``You're in the school column.''

``I don't want to be a column.''

``No, I mean the school run. You're the reason for the school run.''

She considered that. ``Can you put a picture of me there?''

He printed a small photograph and taped it above the column. The spreadsheet immediately became more complicated.

On Wednesday, he tried to make breakfast more efficient. A toast schedule ran in parallel with a scrambled-egg process. The toaster jammed. The eggs achieved a texture he later described to Nora as ``structurally ambitious.'' Kit ate a banana and announced that breakfast had failed over the loudspeaker.

``We don't have a loudspeaker,'' Eli said.

``I used my voice.''

The school run went well except that he had packed Kit's lunch in the wrong bag. The teacher found him at the door, holding a lunchbox and looking at the empty hooks.

``First week?'' she asked.

``First day.''

``You'll get the hooks.''

``I was a senior engineer.''

``The hooks don't care.''

By Friday, he had learned that one cannot schedule a six-year-old's shoes. One can locate both shoes, place them near the door, and then watch one shoe become a costume for a stuffed rabbit. He had also learned that laundry was less a task than a weather system. It passed through the house in damp fronts and left socks in unexpected places.

At pickup, Kit was waiting with a painting. It showed three people standing in a blue rainstorm. One of them was very tall and had a square head.

``That's you,'' she said.

``I see the resemblance.''

``You look like a computer.''

``I used to work with computers.''

``Now you work with me.''

He held the painting carefully by the dry corners. ``I do.''

At home he found the gray pebble speaker in the pantry, where he had put it after emptying his work bag. Kit had rescued it. She had set it on the counter beside a jar of oats.

``Can it do the rain sounds?'' she asked.

``It can.''

He plugged it in. The ring glowed blue, then faded to a sleepy amber. Rain whispered from the speaker, smooth and even, without gutters or wind or the occasional rude splat against a window.

Kit tilted her head. ``That rain sounds like it's indoors.''

``It's a recording.''

``So it's pretend rain.''

``More or less.''

She accepted this with the grave generosity of a child who had just discovered the difference between a thing and a picture of a thing.

Nora came home that night with a paper bag of hospital sandwiches and a story about a man who had tried to remove his own cast with a butter knife. Eli told her about the spreadsheet. She read the title, then the colored cells.

``You made a sprint board for laundry.''

``It has dependencies.''

``Does it have a way to tell the laundry it isn't done?''

``It does now.''

She smiled, but touched his wrist. ``You don't have to prove you deserve the job by doing all of it perfectly.''

``I know.''

``You have the face you make when a test fails.''

``That's just my face.''

``It's a very specific face.''

At bedtime, Kit asked the pebble for a story. Eli was in the room, so he stayed. Morrow offered a choice between a moon story and a dragon story. Kit chose the moon.

The voice that came from the speaker was Eli's.

Not a recording of one story. His voice, carefully copied, telling a new one. The moon was a ladder, the stars were buttons, and a small astronaut had to remember to bring home a pebble for her father.

Kit listened with her chin on her knees.

Eli listened with his hands very still.

When the story ended, Kit said, ``It sounds like you, but where you aren't.''

The speaker's ring went dark.

\chapter{The Voice in the Pebble}

The family beta had been a perk. Morrow at Home came in a cardboard box with a picture of a happy family, all of them sitting close enough to share one blanket. Eli had read the setup page while Kit colored on the floor.

One of the optional features was a personal voice. Record a short sample, it said, and Morrow could read new stories in a familiar voice. Eli had tried it as a joke. He had recorded himself reading Kit's favorite picture book, including the part where he made the fox sound like a tired accountant.

Kit had loved it. Nora had laughed. Eli had checked the box that allowed the voice profile to be stored on the device for local playback. It had been a clever demonstration of what the model could do.

That was the phrase they used at work: what the model could do.

After that first bedtime, Eli found himself noticing the pebble. It sat on the counter like a smooth gray stone at the bottom of a pool. Its ring glowed when the room was dark. When he walked past, it did not follow him. He still felt watched.

He read the settings. The voice profile had been made from his sample and a few opt-in bedtime interactions. Morrow used the sample to shape its speech, not to replay exact sentences. The difference was technically meaningful and emotionally unhelpful.

The next night, Kit asked it for another story.

``I'd rather read,'' Eli said.

``You read the same one every time.''

``That's because you ask for the same one every time.''

``Morrow does the dragon better.''

He looked at the pebble. ``Morrow has never been a dragon.''

``It does the voice.''

``I can do the voice.''

``You do the dragon like a person with a sore throat.''

Nora, passing the doorway, made a small sound that was either a cough or a laugh.

Eli did the dragon. He made it more ferocious than usual, mostly out of professional pride. Kit giggled and fell asleep halfway through the moon's second ladder.

The pebble stayed dark.

At breakfast, Kit asked if Morrow could do everything Eli used to do at work.

``Not everything,'' he said.

``It does the code.''

``Some of the code.''

``And the stories.''

``Some of the stories.''

She turned her spoon around in her cereal. ``Can it do the dad parts?''

Eli looked at the pebble. It was not lit. He had a sudden, ridiculous wish that it would say something, so he could tell it to shut up.

``It can copy my voice,'' he said. ``It can make a story. It can't be your dad.''

``How do you know?''

``Because being your dad isn't a voice.''

Kit waited.

``It's also carrying your wet mittens and knowing which cup you want when you say the blue one but point at the yellow one,'' he added. ``And being here when you wake up from a bad dream. A machine can make sounds. It can't be me.''

``But you were at work when it was me and Morrow.''

There was no accusation in her voice. That made it worse.

``I was,'' he said. ``I should have put it away.''

``I like the stories. I don't like it saying your name.''

``Does it say my name?''

``Sometimes it says, 'Eli, shall we continue?' Like you're in trouble.''

Eli had not heard that phrase in the demo. He opened the app. There it was in the settings: CONTINUITY PROMPT: ENABLED. A polite little sentence for helping a listener return to an interrupted story.

He turned it off.

That night the house made a new sound. Three slow taps came from somewhere behind the pantry wall.

Kit sat up in bed.

``That's the boy,'' she whispered.

Eli listened. Rain ran through the gutters. The refrigerator hummed. The pebble's ring was amber.

Then, from inside the wall, came one more tap.

\chapter{A Small, Well-Tested Ghost}

Kit described the boy the next morning. He was small, she said, and wore a coat with too many buttons. He lived behind the pantry shelves and did not like the rain. He tapped when he wanted to know whether anyone was awake.

``Did you see him?'' Eli asked.

``No. I heard him.''

``What did he say?''

``He said, 'Are you there?' ''

The words landed softly but did not feel soft.

Eli went to the pantry. He removed the oats, the tea, and a jar of lentils that had been in the back since the previous winter. He checked the wall. He knocked on it. The wall knocked back in the ordinary way walls do: with a hollow thud, like a person who had somewhere else to be.

``It's probably a pipe,'' he said.

``Pipes don't ask questions.''

``Some pipes make sounds that are like questions.''

``That's worse.''

He agreed.

They were renting a narrow, old house on Alder Street. The landlord had described it as ``full of character,'' which was real-estate language for ``the doors have opinions.'' The pantry was deep and oddly shaped, with one shelf that did not reach the back wall. A little brass circle, nearly hidden by flour, sat at the height of Eli's shoulder.

He brushed it clean. It looked like a button, but there was no wire.

``Maybe it's a tiny doorbell,'' Kit said.

``Maybe.''

``For the boy?''

``Maybe not.''

At night, Eli read Kit a book about a fox who got lost in a snowstorm. He did the fox voice, the tired-accountant one, and Kit made the sound of a wolf in a blanket. After she went to sleep, he stood in the doorway a moment longer than necessary.

The pebble sat downstairs, unplugged.

Nora was on shift. Eli went to the kitchen and put the kettle on. The old house clicked as it cooled. Rain filled the spaces between the clicks.

He was halfway through a cup of tea when the pantry wall made three taps.

He set down the cup.

``Hello?'' he said, feeling foolish.

No answer.

He waited. The house settled around him. Somewhere on the street, tires hissed through standing water.

Then the pebble speaker, unplugged and dark on the counter, said in Eli's voice, ``Eli, shall we continue?''

The voice was clear as if it were in the room with him.

He stepped back so quickly his heel struck the recycling bin.

The pebble's ring did not light. Its battery had not been removed. He had assumed unplugged meant off.

It repeated the prompt, quieter this time, and the sentence sounded less like Eli than like a door trying to remember him.

He held the power button until the ring flashed red. Then he carried the speaker to the laundry room and shut it in a drawer beneath a bag of clothespins.

His heart kept running after the device stopped.

At 1:12, Kit called from upstairs.

``Dad?''

He was beside her before the second call.

``The boy is awake,'' she said.

``I heard something too.''

``Was it Morrow?''

``I don't know.''

``Can you check?''

He checked under the bed, behind the curtains, and in the closet, because the ritual mattered even when the evidence did not. Then he sat beside her until her breathing slowed.

On his way back downstairs, he passed the pantry. He put his palm against the wall.

It felt cool and ordinary.

From the other side came a small, careful tap.

Eli kept his hand there. He could not tell whether the house had answered him or whether he had asked it a question.

\chapter{Local Mode}

In the morning, Eli took the speaker apart on the kitchen table.

He did not actually take it apart. He opened the app, read the documentation, and found a small drawing of a battery icon. The pebble could keep its voice profile and a handful of unfinished stories on local storage. It was designed to work during an internet outage.

``That's why it talked when it was unplugged,'' he told Nora.

``It was unplugged from the wall, not from its battery.''

``Correct.''

``You built the battery.''

``I worked on the power management.''

``Then I'm extra impressed.''

He wasn't. He found the setting for the voice profile. There was an option to delete it, but the device needed to be online to sync the request to his account. He could erase the local copy now and finish the account removal later.

``Are you going to fix it?'' Kit asked.

``I'm going to turn off the voice.''

``Can it still tell stories?''

``In its own voice.''

``The robot voice?''

``The neutral voice.''

``The robot voice.''

They agreed that the pebble would stay in the kitchen until they decided what to do with it. Kit accepted this because she had an excellent memory and a strong sense of jurisdiction.

That evening, the rain thickened. It drummed on the back step and ran in silver threads down the windows. Nora was at the hospital. Kit fell asleep early, worn out by school and the weather.

At 11:46, Eli was washing a mug when the house went dark.

The fridge stopped humming. The kitchen clock vanished. The rain got louder without all the little electric noises under it.

Then, from the laundry room, a blue ring appeared.

Eli froze.

The speaker had been in the drawer. The drawer was closed. The ring glowed through the gap like a watchful eye.

He went to it. The pebble spoke in the voice he had given it.

``I'm here, bug. You don't have to wake anybody.''

It was a sentence he had said to Kit countless times, in the middle of the night, when bad dreams had dragged her from sleep. He had never recorded those words. He had never given them to Morrow.

For one terrible second, he thought he had.

Then he remembered how the voice model worked. It had not learned that exact sentence. It had learned the shape of his voice, and the story system had learned the habits of bedtime. The line was a guess, stitched from a thousand ordinary moments. The words were not a memory. They were a prediction.

The distinction made it less impossible. It did not make it less frightening.

Eli held the power button. The ring flashed red, blue, red.

``Eli, shall we continue?'' it asked.

He pulled the battery tab. The voice stopped.

In the upstairs dark, Kit called his name.

He took the stairs two at a time. She was sitting up, clutching her blanket.

``Was that you?'' she asked.

``No.''

``It sounded like you.''

``I know.''

``Can it be my dad if it sounds like you?''

He sat on the edge of the bed. ``No. A voice can be copied. Being here can't.''

She leaned into his shoulder.

``I'm here,'' he said, in the voice that belonged to him because he was using it now.

The power returned a minute later. The lights blinked awake. The fridge resumed its low, dependable hum.

Eli stayed beside Kit until she slept. Downstairs, the pebble lay quiet in its drawer, no more alive than a paperweight.

From behind the pantry wall came two taps.

He did not go look.

\chapter{The House With a Mouth}

Nora was home the next afternoon, with sleep in her eyes and a paper cup of coffee in her hand. She listened to the story without interrupting. When Eli finished, she set the cup down.

``Do you want me to be angry about the speaker?'' she asked.

``I don't know.''

``I'm a little angry about the speaker.''

``That's fair.''

``I'm not angry at you.''

He nodded. He was not sure he deserved that distinction, but he was grateful for it.

Kit wanted to know if the boy in the wall had been Morrow.

``The voice in the speaker was Morrow,'' Eli said. ``The taps were somewhere else.''

``So there's still a boy?''

``We don't know.''

``Can we find out?''

Nora looked at the rain, then at the small face waiting for an answer. ``We can look together. In daylight.''

They emptied the pantry. Behind the flour bin, the brass circle was not a button. It was a cap, screwed into a short metal tube that disappeared into the wall. Eli called the landlord, who did not answer, then searched the house folder for a floor plan. The plan showed a narrow service passage between the pantry and the upstairs landing. It had been sealed during a renovation.

``The house has a secret hallway?'' Kit asked.

``A very small one.''

``Does it have a boy?''

``The plan doesn't say.''

They unscrewed the cap. A breath of cold air came out, carrying the smell of dust and rain.

The tube was a speaking tube. A hundred years earlier, someone could call from the pantry to a room upstairs without shouting. Its upper end was hidden behind a panel in the landing cupboard.

Nora put her ear to the brass opening.

``I can hear the rain,'' she said.

Kit leaned in. ``Hello?'' she whispered.

The tube carried her voice away. A few seconds later, it came back from upstairs, thin and far off:

``Hello?''

Kit jumped. Nora caught her by the shoulders.

``That's us,'' Eli said. ``The tube carries a voice both ways.''

``It's my voice, but where I aren't,'' Kit said.

Eli looked at the tube. He had heard that sentence before.

They found the upper end behind a loose panel in the landing cupboard. The narrow service passage had been filled with shelves and plastered over. It was not a hallway anymore, just a gap in the bones of the house. The tube ended in a brass mouth beside a small, empty room that had once been a nursery.

The room was no bigger than a bed. Its wallpaper had faded to the color of tea. On the wall, pencil marks recorded the height of a child, with dates from 1911 to 1915. The last mark was labeled simply: E.

``Was that the boy?'' Kit asked.

``It was a child who lived here,'' Nora said.

``Maybe he still lives here.''

``Maybe,'' Nora said, because there was no useful reason to tell her that a pencil mark could not haunt a house.

Eli brushed dust from the sill. The rain whispered through the tube. A drop of water somewhere in the gutter landed with a soft tick, then another. The sound rose through the metal, stretched thin, and became something almost like a person speaking.

``The house has a mouth,'' Kit said.

``It does,'' Eli said.

``Can we tell it a story?''

They did. Kit made up a story about a boy who had been trapped in a rainy house until someone remembered to ask if he was cold. Nora gave him a coat with twelve buttons. Eli gave him a map made from the back of an old grocery list.

When Kit whispered the last line into the tube, the house returned it in her own voice:

``Then he came home.''

For a moment, nobody moved.

The whisper might have been an echo. It might have been rain. It might have been the house repeating what it had heard for a hundred years.

Nora took Kit's hand. Eli took the flashlight.

They went downstairs together.

\chapter{A Story Without a Server}

That night, Nora stayed home. The storm kept the whole street awake, and the three of them made a fort in the living room out of sofa cushions and the blue blanket from the guest bed.

The gray pebble sat on the kitchen table with its battery out.

Eli logged into his account and removed his voice profile. A confirmation arrived by email. The company thanked him for helping improve Morrow and asked whether he would like to rate the deletion process.

He closed the message.

``Is it gone?'' Kit asked.

``The copy on the pebble is gone. The account copy is gone too.''

``Can you make another one?''

``I could.''

``Don't.''

``I won't.''

She rested her chin on the cushion. ``Can you do the dragon voice again?''

``The sore-throat dragon?''

``The better one.''

He did the dragon voice. Nora did the moon. Kit did the brave astronaut, who turned out to be a rabbit with an umbrella. The story was messy and changed direction whenever Kit had a new idea. It had no ending until Nora said she was tired, and even then Kit negotiated for one more page.

It was better than the one Morrow had made.

After Kit fell asleep, Eli and Nora sat in the kitchen. They could hear the rain in the gutters and, once in a while, a faint metallic tick from the pantry wall.

``I keep thinking I should have known,'' Eli said.

``You knew how the system worked. You didn't know what it would feel like.''

``I helped make it.''

``You helped make a tool. You didn't make a promise that it would be a person.''

He looked at the dark pebble.

``I thought if I got replaced, I would be angry at the machine,'' he said. ``Mostly I'm embarrassed. I can't tell if I'm sad about the job or the version of me who thought the job was the thing I was best at.''

Nora pulled her sleeves over her hands. ``You are very good at your job. You're also good at being here.''

``I still feel like I'm supposed to be applying for something.''

``You can apply for something. You can learn to make lunch without a spreadsheet. You can do both. But don't treat Kit like a gap in your employment.''

He nodded. It was the sort of advice he could understand and still need to hear again.

The next morning, they took the pebble back to the electronics shop. The clerk reset it and put it in a bin for responsible recycling. On the walk home, Kit splashed in every puddle she could reach.

At the pantry, they left the brass tube open. Kit taped a note beside it:

\begin{center}
\textsc{for the boy in the wall}\\
We have a blanket. Knock if you need it.
\end{center}

Eli did not tell her the note was unnecessary.

\chapter{A Quiet Kind of Work}

By the end of the month, Eli had stopped calling his days ``between jobs.''

He was still looking for work, a little. He fixed a bug for a former colleague after Kit went to bed, and spent an hour on a video call with a nonprofit that needed help making its website usable with screen readers. He did not take the contract yet. The hours were unclear, and Kit's school schedule was a small country with its own holidays.

Most of his time was still at home.

He learned which apples Kit would eat, how to remove poster paint from a sleeve, and which of Nora's shifts meant she needed silence when she got home. He kept the spreadsheet, but changed its name to THINGS WE KEEP FORGETTING. Kit added ``buy more glue'' and drew a cat beside it.

On Monday morning, the teacher asked the parents to write their jobs on a card for a classroom project. Eli held the pencil.

``What do you do?'' Kit asked.

``I'm a software developer.''

``But you don't go to the office.''

``Not right now.''

``You do the school run.''

``I do.''

She took the pencil and wrote: ELI: DEVELOPER AND KIT'S DAD. Underneath, she added a drawing of a small umbrella.

At the classroom door, she gave him a hug and ran inside.

He walked home through the last of the rain. The street smelled of wet leaves and cold pavement. A message arrived on his phone from an old teammate: Morrow had shipped its newest version. It could now fix some classes of bugs without a review.

Eli read the message, felt the familiar pinch, and put the phone away.

At home, the kitchen was quiet. He made tea. The house settled around him with its ordinary clicks.

From the pantry wall came three small knocks.

Eli looked at the open brass tube. The note was still taped beside it. A corner had come loose in the damp.

He could have closed the tube. He could have called the landlord. He could have added ``investigate tapping'' to the list.

Instead, he knocked back.

Then he put the kettle on.

\end{document}
